\section{Conclusion}

\begin{frame}{Synthèse de la 2ème Partie}
    Devenir un plongeur autonome Niveau 3 demande une maîtrise globale :
    
    \begin{itemize}
        \item \textbf{La Décompression :} Savoir utiliser son ordinateur, gérer les procédures de secours (tables) et s'adapter aux autres membres de la palanquée.
        \item \textbf{Les Risques :} Prévenir la narcose, l'essoufflement et le froid par une planification adaptée et une connaissance de soi.
        \item \textbf{L'Organisation :} En l'absence de Directeur de Plongée, vous devez tout gérer : du choix du site à la fiche de sécurité et aux secours.
    \end{itemize}
\end{frame}

\begin{frame}{Le mot de la fin}
    \begin{block}{Responsabilité et Autonomie}
        L'autonomie est une liberté qui implique de lourdes responsabilités. Vous êtes désormais garants de votre sécurité et de celle de votre binôme.
    \end{block}

    \vspace{0.5cm}
    \textbf{Avant chaque mise à l'eau :}
    \begin{enumerate}
        \item Vérifiez le matériel de sécurité du bateau.
        \item Accordez-vous sur les paramètres (Prof, Durée, DTR).
        \item Communiquez ! (Signes, procédures).
    \end{enumerate}

    \begin{center}
        \vfill
        \Large \textbf{Bonnes plongées et bonne réussite à l'examen !!!}
    \end{center}
\end{frame}

\begin{frame}
    \centering
    \Huge \textbf{Des questions ?}
\end{frame}